% !TEX program = xelatex
% ============================================================================
%  POSN Computer Qualification — Programming (Python)
%  Topic: การโปรแกรมเบื้องต้น ภาษาไพธอน
% ============================================================================

\documentclass[11pt, aspectratio=169]{beamer}
% ============================================================================
%  POSN Computer Qualification — Shared Beamer Preamble
%  Used by all topic slide decks. Compile with XeLaTeX.
% ============================================================================

% ---------- Thai language & font ----------
\usepackage{iftex}
\usepackage{fontspec}

% XeTeX-specific: enable Thai line breaking
\ifXeTeX
  \XeTeXlinebreaklocale{th}
\fi
% Noto Sans Thai — body text (clean modern sans-serif, handles all Thai)
\setmainfont{Noto Sans Thai}[
  UprightFont = Noto Sans Thai Light,
  BoldFont    = Noto Sans Thai Medium,
  Scale       = 0.9
]
\setsansfont{Noto Sans Thai}[
  UprightFont = Noto Sans Thai Light,
  BoldFont    = Noto Sans Thai Medium,
  Scale       = 0.9
]

% Noto Sans Thai — clean modern sans-serif for structural elements
\newfontfamily\headingfont{Noto Sans Thai}[
  UprightFont = Noto Sans Thai Medium,
  BoldFont    = Noto Sans Thai Bold,
  Scale       = 0.9
]

% --- Heading font for structural elements ---
\setbeamerfont{title}{family=\headingfont, series=\bfseries, size=\huge}
\setbeamerfont{subtitle}{family=\headingfont, series=\mdseries}
\setbeamerfont{author}{family=\headingfont, series=\mdseries}

% --- Custom Title Page ---
\defbeamertemplate{title page}{custom}[1][]{%
  \centering
  \vspace*{2cm}
  % Title — in white
  {\usebeamerfont{title}\color{white}\inserttitle\par}
  \vspace{0.6cm}
  % Subtitle — in white
  {\usebeamerfont{subtitle}\color{white}\insertsubtitle\par}
  \vspace{2.5cm}
  % Author — blue filled rounded box, white text
  \begin{center}
    \begin{tcolorbox}[arc=3mm, colback=boxBlue, colframe=boxBlue!80,
                      width=0.42\textwidth, halign=center,
                      left=1.5em, right=1.5em, top=0.8em, bottom=0.8em,
                      boxrule=1.5pt, coltext=white]
      \usebeamerfont{author}\insertauthor
    \end{tcolorbox}
  \end{center}
}
\setbeamertemplate{title page}[custom]
\setbeamerfont{frametitle}{family=\headingfont, series=\bfseries}
\setbeamerfont{section in toc}{family=\headingfont}
\setbeamerfont{page number in head/foot}{family=\headingfont}
\setbeamerfont{section title}{family=\headingfont, series=\bfseries}
\setbeamerfont{subsection title}{family=\headingfont}

% Monospace font (for code/verbatim)
\setmonofont{Latin Modern Mono}[Scale=0.9]

% ---------- Math ----------
\usepackage{amsmath, amssymb}

% ---------- Graphics ----------
\usepackage{graphicx}
\usepackage{tikz}

% ---------- String parsing (for section headers) ----------
\usepackage{xstring}

% ---------- Colored boxes ----------
\usepackage[skins, breakable, xparse]{tcolorbox}
% Note: we do NOT load the 'most' option — it predefines example/solution environments
% that would conflict with our custom boxes.

% --- Color definitions ---
\definecolor{boxBlue}{RGB}{41, 128, 185}
\definecolor{boxGreen}{RGB}{39, 174, 96}
\definecolor{boxOrange}{RGB}{230, 126, 34}
\definecolor{boxPurple}{RGB}{142, 68, 173}
\definecolor{boxBlueBg}{RGB}{235, 245, 255}
\definecolor{boxGreenBg}{RGB}{235, 255, 240}
\definecolor{boxOrangeBg}{RGB}{255, 245, 235}
\definecolor{boxPurpleBg}{RGB}{248, 240, 255}

% --- Explanation box (blue) — for definitions, concepts, notes ---
\newtcolorbox{explanation}[1][]{
  enhanced,
  breakable,
  arc         = 2mm,
  colback     = boxBlueBg,
  colframe    = boxBlue!50,
  title       = #1,
  fonttitle   = \bfseries\color{white},
  coltitle    = white,
  colbacktitle= boxBlue,
  attach boxed title to top left = {yshift=-2mm, xshift=2mm},
  boxed title style = {arc=2mm, size=small},
  before skip = 8pt,
  after skip  = 8pt
}

% --- Example box (green) — for worked examples ---
\newtcolorbox{exambox}[1][]{
  enhanced,
  breakable,
  arc         = 2mm,
  colback     = boxGreenBg,
  colframe    = boxGreen!50,
  title       = #1,
  fonttitle   = \bfseries\color{white},
  coltitle    = white,
  colbacktitle= boxGreen,
  attach boxed title to top left = {yshift=-2mm, xshift=2mm},
  boxed title style = {arc=2mm, size=small},
  before skip = 8pt,
  after skip  = 8pt
}

% --- Solution box (orange) — for solutions and answers ---
\newtcolorbox{solbox}[1][]{
  enhanced,
  breakable,
  arc         = 2mm,
  colback     = boxOrangeBg,
  colframe    = boxOrange!50,
  title       = #1,
  fonttitle   = \bfseries\color{white},
  coltitle    = white,
  colbacktitle= boxOrange,
  attach boxed title to top left = {yshift=-2mm, xshift=2mm},
  boxed title style = {arc=2mm, size=small},
  before skip = 8pt,
  after skip  = 8pt
}

% --- Intuition box (purple) — for plain-language explanations, analogies ---
\newtcolorbox{intuition}[1][]{
  enhanced,
  breakable,
  arc         = 2mm,
  colback     = boxPurpleBg,
  colframe    = boxPurple!50,
  title       = #1,
  fonttitle   = \bfseries\color{white},
  coltitle    = white,
  colbacktitle= boxPurple,
  attach boxed title to top left = {yshift=-2mm, xshift=2mm},
  boxed title style = {arc=2mm, size=small},
  before skip = 8pt,
  after skip  = 8pt
}

% ---------- Beamer theme ----------
\usetheme{Madrid}
\usecolortheme{default}

% --- Custom beamer colors (blue accent matching explanation boxes) ---
\setbeamercolor{structure}{fg=boxBlue}
\setbeamercolor{palette primary}{fg=white, bg=boxBlue}
\setbeamercolor{palette secondary}{fg=white, bg=boxBlue!80}
\setbeamercolor{palette tertiary}{fg=white, bg=boxBlue!60}
\setbeamercolor{block title}{fg=white, bg=boxBlue}
\setbeamercolor{block body}{fg=black, bg=boxBlueBg}

% Remove navigation symbols for clean look
\setbeamertemplate{navigation symbols}{}

% Note: Frame title prefix removed — \addtobeamertemplate conflicts with beamer internals

% --- Outline (Table of Contents) styling ---
\setbeamerfont{section in toc}{family=\headingfont, series=\mdseries}
\setbeamerfont{subsection in toc}{family=\headingfont, series=\mdseries}
\setbeamertemplate{section in toc}{%
  \leavevmode\leftskip=1.5em
  \textcolor{boxBlue}{\textbf{\large\inserttocsectionnumber.}}%
  \quad\inserttocsection\par
  \vspace{0.2em}
}
\setbeamertemplate{subsection in toc}{%
  \leavevmode\leftskip=4em
  \textcolor{boxBlue!60}{\small$\bullet$}%
  \quad\inserttocsubsection\par
  \vspace{0.1em}
}

% Section header slides — Thai in white, English in smaller light blue below
\AtBeginSection{
  \begin{frame}[plain]{}
    \vspace*{2cm}
    \begin{beamercolorbox}[sep=12pt, center, rounded=true, shadow=true]{palette primary}
      \IfSubStr{\secname}{(}{%
        \StrBefore{\secname}{(}[\sectionthai]%
        \StrBetween[1,1]{\secname}{(}{)}[\sectioneng]%
        \begin{tabular}{@{}c@{}}
          {\usebeamerfont{title}\sectionthai} \\[0.2em]
          {\usebeamerfont{subtitle}\color{boxBlue!60}\sectioneng}
        \end{tabular}%
      }{%
        {\usebeamerfont{title}\secname\par}%
      }
    \end{beamercolorbox}
    \vfill
  \end{frame}
}

% Subsection header slides — smaller, centered
\AtBeginSubsection{
  \begin{frame}[plain]{}
    \vfill
    \centering
    {\usebeamerfont{section title}\insertsubsectionhead\par}
    \vfill
  \end{frame}
}

% Minimal footer: page number only, right-aligned
\setbeamertemplate{footline}{
  \hfill\usebeamerfont{page number in head/foot}\insertframenumber{} / \inserttotalframenumber\hspace*{1em}\vspace*{0.5ex}
}

% ---------- Spacing & readability ----------
\linespread{1.2}

% ---------- Useful commands ----------

% Highlight important text in blue
\newcommand{\highlight}[1]{\textcolor{boxBlue}{\textbf{#1}}}

% Math set notation helpers
\newcommand{\R}{\mathbb{R}}
\newcommand{\N}{\mathbb{N}}
\newcommand{\Z}{\mathbb{Z}}
\newcommand{\Q}{\mathbb{Q}}

% ============================================================================
%  End of preamble
% ============================================================================


% --- Code font: Courier New ---
\setmonofont{Courier New}[
  Path       = ../../fonts/CourierNew/,
  Extension  = .ttf,
  UprightFont = cour,
  Scale       = 0.85
]

\newcommand{\py}[1]{\texttt{#1}}

\title[การโปรแกรม]{\textcolor{white}{การโปรแกรมเบื้องต้น} \textcolor{boxBlue!60}{(Python)}}
\subtitle{Data Types • Operators • Conditionals • While Loop • List • String }
\author[None W. (Yuan)]{None W. (Yuan)}
\date{}

\begin{document}

\begin{frame}[plain]\titlepage\end{frame}
\begin{frame}{สารบัญ (Outline)}\tableofcontents\end{frame}

\begin{frame}{ก่อนเรียน (Prerequisites)}
  \begin{explanation}[สิ่งที่ควรรู้ก่อนเรียน]
    \begin{itemize}
      \item \textbf{คณิตศาสตร์พื้นฐาน} — การบวก ลบ คูณ หาร, การเปรียบเทียบ, ตรรกศาสตร์
      \item \textbf{ตัวแปร} — แนวคิดเรื่องการเก็บค่า (เหมือนคณิตศาสตร์)
      \item \textbf{คอมพิวเตอร์} — การใช้งานเครื่องคอมพิวเตอร์พื้นฐาน
    \end{itemize}
  \end{explanation}
  \begin{intuition}[การเขียนโปรแกรมคืออะไร?]
    การเขียนโปรแกรมคือการ \textbf{สั่งให้คอมพิวเตอร์ทำงานทีละขั้นตอน}
    โดยใช้ภาษาที่คอมพิวเตอร์เข้าใจ — ในที่นี้คือ \highlight{ภาษาไพธอน (Python)}
  \end{intuition}
\end{frame}

% ===========================================================================
%  SECTION 1: ชนิดข้อมูลและตัวแปร
% ===========================================================================
\section{ชนิดข้อมูลและตัวแปร (Data Types \& Variables)}

\begin{frame}{ชนิดข้อมูลพื้นฐาน (Basic Data Types)}
  \begin{explanation}[4 ชนิดข้อมูลสำคัญ]
    \begin{columns}[T]
      \column{0.5\textwidth}
        \begin{itemize}
          \item \highlight{\py{int}} — จำนวนเต็ม (integer)
          \begin{itemize}
            \item เช่น \py{5}, \py{-3}, \py{0}, \py{100}
          \end{itemize}
          \item \highlight{\py{float}} — จำนวนจริง/ทศนิยม
          \begin{itemize}
            \item เช่น \py{3.14}, \py{-0.5}, \py{2.0}
          \end{itemize}
        \end{itemize}
      \column{0.5\textwidth}
        \begin{itemize}
          \item \highlight{\py{bool}} — ค่าความจริง (Boolean)
          \begin{itemize}
            \item \py{True}, \py{False}
          \end{itemize}
          \item \highlight{\py{str}} — ข้อความ (string)
          \begin{itemize}
            \item เช่น \py{"Hello"}, \py{"123"}
          \end{itemize}
        \end{itemize}
    \end{columns}
  \end{explanation}
  \begin{exambox}[ตัวอย่าง: การเช็คชนิดข้อมูล (ไม่ออกสอบ)]
    \py{type(5)} $\rightarrow$ \py{<class 'int'>} \quad
    \py{type(3.14)} $\rightarrow$ \py{<class 'float'>}

    \py{type(True)} $\rightarrow$ \py{<class 'bool'>} \quad
    \py{type("hello")} $\rightarrow$ \py{<class 'str'>}
  \end{exambox}
\end{frame}

\begin{frame}{ตัวแปรและการกำหนดค่า (Variables \& Assignment)}
  \begin{explanation}[การกำหนดค่า]
    \textbf{ตัวแปร} คือชื่อที่ใช้อ้างถึงค่าที่เก็บในหน่วยความจำ

    ใช้เครื่องหมาย \py{=} ในการ \highlight{กำหนดค่า}: \py{x = 5} หมายถึง เก็บค่า $5$ ในตัวแปร \py{x}
  \end{explanation}
  \begin{columns}[T]
    \column{0.45\textwidth}
      \begin{exambox}[ตัวอย่าง: การกำหนดค่า]
        \py{x = 5}

        \py{y = 3}

        \py{z = x + y} \qquad (ได้ \py{z = 8})

        \py{name = "Yuan"}

        \py{is\_valid = True}
      \end{exambox}
    \column{0.55\textwidth}
      \begin{explanation}[กฎการตั้งชื่อตัวแปร]
        \begin{itemize}
          \item ขึ้นต้นด้วยตัวอักษร หรือ \texttt{\_} (underscore)
          \item ประกอบด้วยตัวอักษร ตัวเลข และ \texttt{\_}
          \item \textbf{ห้าม} ขึ้นต้นด้วยตัวเลข
          \item \textbf{ห้าม} ใช้คำสงวน เช่น \py{if}, \py{while}
          \item \textbf{Case-sensitive}: \py{x} กับ \py{X} เป็นคนละตัว
        \end{itemize}
      \end{explanation}
  \end{columns}
\end{frame}

% ===========================================================================
%  SECTION 2: ตัวดำเนินการ
% ===========================================================================
\section{ตัวดำเนินการ (Operators)}

\begin{frame}{ตัวดำเนินการทางคณิตศาสตร์ (Arithmetic Operators)}
  \begin{columns}[T]
    \column{0.52\textwidth}
      \begin{explanation}[เครื่องหมายพื้นฐาน]
        \begin{tabular}{c|l|l}
          \textbf{เครื่องหมาย} & \textbf{ความหมาย} & \textbf{ตัวอย่าง} \\
          \hline
          \py{+}  & บวก     & \py{5 + 3} $\rightarrow$ \py{8} \\
          \py{-}  & ลบ      & \py{5 - 3} $\rightarrow$ \py{2} \\
          \py{*}  & คูณ     & \py{5 * 3} $\rightarrow$ \py{15} \\
          \py{/}  & หาร      & \py{5 / 2} $\rightarrow$ \py{2.5} \\
          \py{//} & หารปัดเศษ & \py{5 // 2} $\rightarrow$ \py{2} \\
          \py{\%} & \textbf{มอดุโล (mod)} & \py{5 \% 2} $\rightarrow$ \py{1} \\
          \py{**} & ยกกำลัง  & \py{2 ** 3} $\rightarrow$ \py{8} \\
        \end{tabular}
      \end{explanation}
    \column{0.5\textwidth}
      \begin{exambox}[ตัวอย่าง: mod (\%)]
        \py{10 \% 3 = 1} \quad ($10 = 3\times3 + 1$)

        \py{17 \% 5 = 2} \quad ($17 = 5\times3 + 2$)

        \py{8 \% 2 = 0} \quad ($8 = 2\times4 + 0$)\\

        \textbf{การใช้งาน:}
        \begin{itemize}
          \item เช็คเลขคู่/คี่: \py{n \% 2 == 0}
          \item หลักหน่วย: \py{n \% 10}
          \item การวนรอบ (cycle)
        \end{itemize}
      \end{exambox}
  \end{columns}
\end{frame}

\begin{frame}{ตัวดำเนินการเปรียบเทียบและตรรกะ}
  \begin{columns}[T]
    \column{0.45\textwidth}
      \begin{explanation}[Comparison]
        \py{==} (เท่ากับ), \py{!=} (ไม่เท่ากับ), \py{<}, \py{>}, \py{<=}, \py{>=}
        \medskip

        \py{5 == 5} $\rightarrow$ \py{True}

        \py{3 != 5} $\rightarrow$ \py{True}
      \end{explanation}
    \column{0.55\textwidth}
      \begin{explanation}[Logical]
        \py{and}, \py{or}, \py{not}
        \medskip

        \py{(5 > 3) and (2 < 4)} $\rightarrow$ \py{True}

        \py{(5 > 3) or (2 > 4)} $\rightarrow$ \py{True}

        \py{not (5 == 5)} $\rightarrow$ \py{False}
      \end{explanation}
  \end{columns}
\end{frame}

\begin{frame}{ลำดับความสำคัญของตัวดำเนินการ (Operator Precedence)}
  \begin{explanation}[ลำดับจากสูงไปตํ่า (ทำก่อน $\rightarrow$ ทำหลัง)]
    \begin{enumerate}
      \item \py{()} — วงเล็บ (สำคัญที่สุด)
      \item \py{**} — ยกกำลัง (ทำจาก\textbf{ขวาไปซ้าย})
      \item \py{*}, \py{/}, \py{//}, \py{\%} — คูณ หาร มอดุโล
      \item \py{+}, \py{-} — บวก ลบ
      \item \py{==}, \py{!=}, \py{<}, \py{>}, \py{<=}, \py{>=} — เปรียบเทียบ
      \item \py{not} — นิเสธ
      \item \py{and} — และ
      \item \py{or} — หรือ (สำคัญน้อยที่สุด)
    \end{enumerate}
  \end{explanation}
\end{frame}

\begin{frame}{ลำดับความสำคัญของตัวดำเนินการ (Operator Precedence)}
  \begin{exambox}[ตัวอย่าง]
    \py{3 + 4 * 2 = 11} (คูณก่อนบวก) \qquad
    \py{(3 + 4) * 2 = 14} (วงเล็บก่อน)

    \py{10 / 2 ** 2 = 2.5} (ยกกำลังก่อนหาร) \qquad
    \py{2 ** 3 ** 2 = 512} (ขวาไปซ้าย)
  \end{exambox}
\end{frame}

% ===========================================================================
%  SECTION 3: การแสดงผล
% ===========================================================================
\section{การแสดงผล (Output)}

\begin{frame}{คำสั่ง \py{print()}}
  \begin{explanation}[ใช้แสดงผลลัพธ์ออกทางหน้าจอ]
    \py{print()} — แสดงค่าแล้วขึ้นบรรทัดใหม่
  \end{explanation}
  \begin{exambox}[ตัวอย่าง: การใช้ print()]
    \py{print("Hello World")} \quad $\rightarrow$ \py{Hello World}

    \py{print(3 + 5)} \quad $\rightarrow$ \py{8}

    \py{x = 10}

    \py{print(x)} \quad $\rightarrow$ \py{10}

    \py{print("x =", x)} \quad $\rightarrow$ \texttt{x = 10}
  \end{exambox}
\end{frame}

% ===========================================================================
%  SECTION 4: เงื่อนไข
% ===========================================================================
\section{เงื่อนไข (Conditionals)}

\begin{frame}{คำสั่ง \py{if} / \py{elif} / \py{else}}
  \begin{explanation}[การตัดสินใจในโปรแกรม]
    \highlight{\py{if} เงื่อนไข:} ทำเมื่อเงื่อนไขเป็นจริง

    \highlight{\py{elif} เงื่อนไข:} เงื่อนไขรอง (else if) เมื่อเงื่อนไขข้างบนไม่เป็นจริง มีหลาย \py{elif} ได้

    \highlight{\py{else}:} ทำเมื่อไม่มีเงื่อนไขใดเป็นจริงเลย
  \end{explanation}

  \begin{intuition}[สำคัญ!]
    \py{if-elif-else} เป็นกลุ่มเดียวกันจะเข้าเงื่อนไขได้แค่เงื่อนไขเดียวในกลุ่ม หลังจากเข้าเงื่อนไขด้านบนแล้วจะไม่สนใจเงื่อนไขด้านล่างอีก
    
    ถ้าจะสร้างกลุ่มใหม่ต้องเริ่ม \py{if} ใหม่
  \end{intuition}
\end{frame}

\begin{frame}{ตัวอย่างและการเยื้อง}
  \begin{columns}[T]
    \column{0.5\textwidth}
      \begin{exambox}[ตัวอย่าง: เกรด]
        \py{score = 85}

        \py{if score >= 80:}

        \qquad \py{print("A")}

        \py{elif score >= 70:}

        \qquad \py{print("B")}

        \py{elif score >= 60:}

        \qquad \py{print("C")}

        \py{else:}

        \qquad \py{print("F")}

        \qquad \py{print("You failed!")}
      \end{exambox}
    \column{0.5\textwidth}
      \begin{explanation}[การเยื้อง (Indentation)]
        \highlight{การเยื้อง (indent) มีความหมาย!}

        ในภาษา Python \textbf{การเว้นวรรคหน้าคำสั่ง}
        เป็นตัวกำหนดว่าคำสั่งนั้นอยู่ในบล็อกใด

        \medskip
        \textbf{กฎ:}
        \begin{itemize}
          \item ใช้ 4 ช่องว่าง (space) หรือ 1 tab
          \item \textbf{ห้ามปนกัน} ในบล็อกเดียวกัน
          \item คำสั่งในบล็อกเดียวกันต้องเยื้องเท่ากัน
        \end{itemize}
      \end{explanation}
  \end{columns}
\end{frame}

\begin{frame}{ตัวอย่าง: การตามรอยเงื่อนไข}
  \begin{exambox}[โจทย์]
    ถ้า \py{x = 5} และ \py{y = 12} ผลลัพธ์ของโปรแกรมเป็นเท่าใด?
  \end{exambox}
  \begin{columns}[T]
    \column{0.45\textwidth}
      \begin{exambox}[โปรแกรม]
        \py{x = 5}

        \py{y = 12}

        \py{if x > 10:}

        \qquad \py{print("A")}

        \py{elif y > 10:}

        \qquad \py{print("B")}

        \py{else:}

        \qquad \py{print("C")}
      \end{exambox}
    \column{0.55\textwidth}
      \begin{solbox}[วิธีตามรอย]
        \py{x = 5}: \py{x > 10} เป็น \py{False} $\rightarrow$ ไม่ทำ

        \py{y = 12}: \py{y > 10} เป็น \py{True} $\rightarrow$ \highlight{ทำ \py{elif}}

        แสดงผล \py{"B"} \quad (ข้าม \py{else})
      \end{solbox}
  \end{columns}
\end{frame}

\begin{frame}{ตัวอย่าง: การตามรอยเงื่อนไข}
  \begin{exambox}[โจทย์]
    ถ้า \py{x = 11} และ \py{y = 12} ผลลัพธ์ของโปรแกรมเป็นเท่าใด?
  \end{exambox}
  \begin{columns}[T]
    \column{0.45\textwidth}
      \begin{exambox}[โปรแกรม]
        \py{x = 11}

        \py{y = 12}

        \py{if x > 10:}

        \qquad \py{print("A")}

        \py{elif y > 10:}

        \qquad \py{print("B")}

        \py{if x > 5:}

        \qquad \py{print("C")}
      \end{exambox}
    \column{0.55\textwidth}
      \begin{solbox}[วิธีตามรอย]
        \py{x = 11}: \py{x > 10} เป็น \py{True} $\rightarrow$ ทำ \py{if} ด้านบน

        แสดงผล \py{"A"} \quad (ข้าม \py{elif})

        \py{x = 11}: \py{x > 5} เป็น \py{True} $\rightarrow$ ทำ \py{if} ด้านล่าง

        แสดงผล \py{"C"}
      \end{solbox}
  \end{columns}
\end{frame}

% ===========================================================================
%  SECTION 5: การวนซํ้า
% ===========================================================================
\section{การวนซํ้า (While Loop)}

\begin{frame}{คำสั่ง \py{while}}
  \begin{explanation}[การทำซํ้าตราบเท่าที่เงื่อนไขเป็นจริง]
    \highlight{\py{while} เงื่อนไข:} ทำคำสั่งในบล็อกซํ้าไปเรื่อย ๆ \textbf{ตราบเท่าที่} เงื่อนไขยังเป็นจริง
  \end{explanation}
  \begin{columns}[T]
    \column{0.45\textwidth}
      \begin{exambox}[ตัวอย่าง: นับ 1 ถึง 5]
        \py{i = 1}

        \py{while i <= 5:}

        \qquad \py{print(i)}

        \qquad \py{i = i + 1}
      \end{exambox}
      \textbf{ผลลัพธ์:} \py{1 2 3 4 5}
    \column{0.55\textwidth}
      \begin{explanation}[การทำงาน]
        \begin{tabular}{c|c|c}
          \textbf{รอบ} & \py{i} & \py{i <= 5?} \\
          \hline
          1 & 1 & True $\rightarrow$ print 1 \\
          2 & 2 & True $\rightarrow$ print 2 \\
          3 & 3 & True $\rightarrow$ print 3 \\
          4 & 4 & True $\rightarrow$ print 4 \\
          5 & 5 & True $\rightarrow$ print 5 \\
          6 & 6 & False $\rightarrow$ \textbf{ออกจากวน} \\
        \end{tabular}
      \end{explanation}
  \end{columns}
\end{frame}

\begin{frame}{\py{break} และ \py{continue}}
  \begin{explanation}[คำสั่งควบคุมการวนซํ้า]
    \highlight{\py{break}} — \textbf{ออกจาก} การวนซํ้าทันที (หยุดเลย)

    \highlight{\py{continue}} — \textbf{ข้าม} คำสั่งที่เหลือ แล้วขึ้นรอบถัดไป
  \end{explanation}
  \begin{columns}[T]
    \column{0.5\textwidth}
      \begin{exambox}[ตัวอย่าง: break]
        \py{i = 1}

        \py{while True:}

        \qquad \py{print(i)}

        \qquad \py{if i == 3: break}

        \qquad \py{i = i + 1}
      \end{exambox}
      \textbf{ผลลัพธ์:} \py{1 2 3}
    \column{0.5\textwidth}
      \begin{exambox}[ตัวอย่าง: continue]
        \py{i = 0}

        \py{while i < 5:}

        \qquad \py{i = i + 1}

        \qquad \py{if i == 3: continue}

        \qquad \py{print(i)}
      \end{exambox}
      \textbf{ผลลัพธ์:} \py{1 2 4 5}
  \end{columns}
\end{frame}

\begin{frame}{ตัวอย่าง: การตามรอย while loop}
  \begin{exambox}[โจทย์]
    ถ้าเริ่มต้น \py{n = 10} ผลลัพธ์ของโปรแกรมเป็นเท่าใด?
  \end{exambox}
  \begin{columns}[T]
    \column{0.4\textwidth}
      \begin{exambox}[โปรแกรม]
        \py{n = 10}

        \py{s = 0}

        \py{while n > 0:}

        \qquad \py{s = s + (n \% 10)}

        \qquad \py{n = n // 10}

        \py{print(s)}
      \end{exambox}
    \column{0.6\textwidth}
      \begin{solbox}[วิธีตามรอย]
        \begin{tabular}{c|c|c|c}
          \textbf{รอบ} & \py{n} & \py{n\%10} & \py{s} \\
          \hline
          เริ่ม & 10 & -- & 0 \\
          1 & 10 & 0 & 0 \\
          2 & 1 & 1 & 1 \\
        \end{tabular}

        \highlight{แสดงผล \py{1}} (ผลรวมเลขโดดของ 10 คือ $1+0=1$)
      \end{solbox}
  \end{columns}
\end{frame}

\begin{frame}{ตัวอย่าง: while loop (ต่อ)}
  \begin{exambox}[โจทย์]
    ถ้า \py{x = 2} และ \py{y = 5} ผลลัพธ์ของโปรแกรมเป็นเท่าใด?
  \end{exambox}
  \begin{columns}[T]
    \column{0.4\textwidth}
      \begin{exambox}[โปรแกรม]
        \py{x = 2}

        \py{y = 5}

        \py{c = 0}

        \py{while x < y:}

        \qquad \py{x = x + 1}

        \qquad \py{y = y - 1}

        \qquad \py{c = c + 1}

        \py{print(c)}
      \end{exambox}
    \column{0.6\textwidth}
      \begin{solbox}[วิธีตามรอย]
        \begin{tabular}{c|c|c|c}
          \textbf{รอบ} & \py{x} & \py{y} & \py{c} \\
          \hline
          เริ่ม & 2 & 5 & 0 \\
          1 & 3 & 4 & 1 \\
          2 & 4 & 3 & 2 \\
        \end{tabular}

        รอบ 3: \py{x = 4, y = 3} $\rightarrow$ \highlight{\py{x < y} = False} $\rightarrow$ จบ

        \highlight{แสดงผล \py{2}}
      \end{solbox}
  \end{columns}
\end{frame}

\begin{frame}{Nested While Loop (while ซ้อน while)}
  \begin{explanation}[แนวคิด]
    \textbf{while loop ซ้อนกัน}: loop ในอยู่ภายใน loop นอก

    loop นอกวน 1 รอบ $\rightarrow$ loop ในวนจนครบ $\rightarrow$ loop นอกวนรอบถัดไป
  \end{explanation}

  \begin{columns}[T]
    \column{0.55\textwidth}
      \begin{exambox}[ตัวอย่าง: สูตรคูณ]
        \py{i = 1}

        \py{while i <= 3:}

        \qquad \py{j = 1}

        \qquad \py{while j <= 3:}

        \qquad\quad \py{print(i, "x", j, "=", i*j)}

        \qquad\quad \py{j = j + 1}

        \qquad \py{i = i + 1}
      \end{exambox}

      ผลลัพธ์: \py{1x1=1}, \py{1x2=2}, \dots, \py{3x3=9} (9 บรรทัด)

    \column{0.45\textwidth}
      \begin{explanation}[การทำงาน]
        \py{i=1}: \py{j=1,2,3} $\rightarrow$ 1x1, 1x2, 1x3

        \py{i=2}: \py{j=1,2,3} $\rightarrow$ 2x1, 2x2, 2x3

        \py{i=3}: \py{j=1,2,3} $\rightarrow$ 3x1, 3x2, 3x3

        \highlight{รวม $3 \times 3 = 9$ รอบ!}
      \end{explanation}
  \end{columns}
\end{frame}

\begin{frame}{ตัวอย่าง: Nested While Loop}
  \begin{exambox}[โจทย์]
    ถ้า \py{a = 2} และ \py{b = 3} ผลลัพธ์ของโปรแกรมเป็นเท่าใด?
  \end{exambox}

  \begin{columns}[T]
    \column{0.5\textwidth}
      \begin{exambox}
        \py{s = 0}

        \py{i = 1}

        \py{while i <= 2:}

        \qquad \py{j = 1}

        \qquad \py{while j <= 3:}

        \qquad\quad \py{s = s + i * j}

        \qquad\quad \py{j = j + 1}

        \qquad \py{i = i + 1}

        \py{print(s)}
      \end{exambox}
    \column{0.5\textwidth}
      \begin{solbox}[วิธีตามรอย]
        \py{i=1}: \py{j=1} $\rightarrow$ \py{s = 0 + 1}

        \qquad\; \py{j=2} $\rightarrow$ \py{s = 1 + 2 = 3}

        \qquad\; \py{j=3} $\rightarrow$ \py{s = 3 + 3 = 6}

        \py{i=2}: \py{j=1} $\rightarrow$ \py{s = 6 + 2 = 8}

        \qquad\; \py{j=2} $\rightarrow$ \py{s = 8 + 4 = 12}

        \qquad\; \py{j=3} $\rightarrow$ \py{s = 12 + 6 = 18}

        \highlight{แสดงผล \py{18}}
      \end{solbox}
  \end{columns}
\end{frame}

% ===========================================================================
%  SECTION 6: แถวลำดับ
% ===========================================================================
\section{แถวลำดับ (1D Array / List)}

\begin{frame}{List — แถวลำดับ 1 มิติ}
  \begin{explanation}[นิยาม]
    \py{list} คือโครงสร้างข้อมูลที่เก็บค่าหลาย ๆ ค่าเรียงต่อกันในตัวแปรเดียว

    สร้างด้วย \py{[ ]} คั่นด้วย \py{,} \quad เช่น \py{arr = [10, 20, 30, 40]}
  \end{explanation}
  \begin{columns}[T]
    \column{0.5\textwidth}
      \begin{explanation}[การเข้าถึงสมาชิก (Indexing)]
        \textbf{ดัชนี (index)} เริ่มจาก \highlight{\py{0}}

        \begin{tabular}{c|c|c|c|c}
          ค่า   & 10 & 20 & 30 & 40 \\
          \hline
          index & \py{0} & \py{1} & \py{2} & \py{3} \\
        \end{tabular}

        \py{arr[0]} $\rightarrow$ \py{10} \qquad
        \py{arr[2]} $\rightarrow$ \py{30}
      \end{explanation}
    \column{0.5\textwidth}
      \begin{exambox}[ตัวอย่าง: list]
        \py{arr = [10, 20, 30, 40]}

        \py{print(arr[0])} \ \ {\small \# 10}

        \py{arr[2] = 99}

        \py{print(arr)} \ \ {\small \# [10, 20, 99, 40]}

        \py{print(len(arr))} \ \ {\small \# 4}
      \end{exambox}
  \end{columns}
\end{frame}

\begin{frame}{ตัวอย่าง: การวน loop กับ list}
  \begin{exambox}[โจทย์]
    ถ้า \py{arr = [3, 7, 2, 8, 5]} ผลลัพธ์ของโปรแกรมเป็นเท่าใด?
  \end{exambox}
  \begin{columns}[T]
    \column{0.45\textwidth}
      \begin{exambox}[โปรแกรม]
        \py{arr = [3, 7, 2, 8, 5]}

        \py{s = 0}

        \py{i = 0}

        \py{while i < len(arr):}

        \qquad \py{if arr[i] \% 2 == 1:}

        \qquad\quad \py{s = s + arr[i]}

        \qquad \py{i = i + 1}

        \py{print(s)}
      \end{exambox}
    \column{0.55\textwidth}
      \begin{solbox}[วิธีตามรอย]
        \py{arr[0] = 3}: \py{3\%2 = 1} (คี่) $\rightarrow$ \py{s = 3}

        \py{arr[1] = 7}: คี่ $\rightarrow$ \py{s = 10}

        \py{arr[2] = 2}: คู่ $\rightarrow$ ข้าม

        \py{arr[3] = 8}: คู่ $\rightarrow$ ข้าม

        \py{arr[4] = 5}: คี่ $\rightarrow$ \py{s = 15}

        \highlight{แสดงผล \py{15}}
      \end{solbox}
  \end{columns}
\end{frame}

% ===========================================================================
%  SECTION 7: สตริง
% ===========================================================================
\section{สตริง (String)}

\begin{frame}{String — ข้อความ}
  \begin{explanation}[นิยาม]
    \py{str} คือข้อมูลประเภทข้อความ สร้างด้วย \py{" "} หรือ \py{' '}

    String ก็คือ \highlight{list ของตัวอักษร} — ใช้ index ได้เหมือน list
  \end{explanation}
  \begin{columns}[T]
    \column{0.5\textwidth}
      \begin{exambox}[ตัวอย่าง: string]
        \py{s = "PYTHON"}

        \py{print(s[0])} \ \ {\small \# P}

        \py{print(s[2])} \ \ {\small \# T}

        \py{print(len(s))} \ \ {\small \# 6}

        \py{a = "Hello"; b = "World"}

        \py{print(a + " " + b)} \ \ {\small \# Hello World}
      \end{exambox}
    \column{0.5\textwidth}
      \begin{explanation}[การดำเนินการกับ string]
        \begin{itemize}
          \item \py{len(s)} — ความยาว
          \item \py{s[i]} — ตัวอักษรที่ index \py{i}
          \item \py{s1 + s2} — ต่อข้อความ
          \item \py{s * 3} — ทำซํ้า \\
          (\py{"Hi"*3} $\rightarrow$ \py{"HiHiHi"})
        \end{itemize}
      \end{explanation}
  \end{columns}
\end{frame}

\begin{frame}{ตัวอย่าง: การวน loop กับ string}
  \begin{exambox}[โจทย์]
    ถ้า \py{s = "COMPUTER"} ผลลัพธ์ของโปรแกรมเป็นเท่าใด?
  \end{exambox}
  \begin{columns}[T]
    \column{0.5\textwidth}
      \begin{exambox}[โปรแกรม]
        \py{s = "COMPUTER"}

        \py{c = 0}

        \py{i = 0}

        \py{while i < len(s):}

        \qquad \py{if s[i] == 'U': break}

        \qquad \py{c = c + 1}

        \qquad \py{i = i + 1}

        \py{print(c)}
      \end{exambox}
    \column{0.5\textwidth}
      \begin{solbox}[วิธีตามรอย]
        \py{s[0] = 'C'} $\rightarrow$ \py{c = 1}

        \py{s[1] = 'O'} $\rightarrow$ \py{c = 2}

        \py{s[2] = 'M'} $\rightarrow$ \py{c = 3}

        \py{s[3] = 'P'} $\rightarrow$ \py{c = 4}

        \py{s[4] = 'U'} $\rightarrow$ \highlight{เจอ 'U', break!}

        \highlight{แสดงผล \py{4}} (จำนวนตัวอักษรก่อนถึง U)
      \end{solbox}
  \end{columns}
\end{frame}

% ===========================================================================
%  SECTION 8: ตัวอย่างประยุกต์
% ===========================================================================
\section{ตัวอย่างประยุกต์}

\begin{frame}{ตัวอย่าง: การนับเลขโดด}
  \begin{exambox}[โจทย์]
    ถ้า \py{n = 12345} ผลลัพธ์ของโปรแกรมเป็นเท่าใด?
  \end{exambox}
  \begin{columns}[T]
    \column{0.4\textwidth}
      \begin{exambox}[โปรแกรม]
        \py{n = 12345}

        \py{c = 0}

        \py{while n > 0:}

        \qquad \py{c = c + 1}

        \qquad \py{n = n // 10}

        \py{print(c)}
      \end{exambox}
    \column{0.6\textwidth}
      \begin{solbox}[วิธีตามรอย]
        \begin{tabular}{c|c|c}
          \textbf{รอบ} & \py{n} & \py{c} \\
          \hline
          เริ่ม & 12345 & 0 \\
          1 & 1234 & 1 \\
          2 & 123 & 2 \\
          3 & 12 & 3 \\
          4 & 1 & 4 \\
          5 & 0 & 5 \\
        \end{tabular}

        \highlight{แสดงผล \py{5}} (จำนวนเลขโดด)
      \end{solbox}
  \end{columns}
\end{frame}

\begin{frame}{ตัวอย่าง: การหาเลขคู่}
  \begin{exambox}[โจทย์]
    ถ้า \py{arr = [4, 7, 1, 8, 3]} ผลลัพธ์ของโปรแกรมเป็นเท่าใด?
  \end{exambox}
  \begin{columns}[T]
    \column{0.5\textwidth}
      \begin{exambox}[โปรแกรม]
        \py{arr = [4, 7, 1, 8, 3]}

        \py{c = 0}

        \py{i = 0}

        \py{while i < 5:}

        \qquad \py{if arr[i] \% 2 == 0:}

        \qquad\quad \py{c = c + 1}

        \qquad \py{i = i + 1}

        \py{print(c)}
      \end{exambox}
    \column{0.5\textwidth}
      \begin{solbox}[วิธีตามรอย]
        \py{4 \% 2 = 0} $\rightarrow$ \py{c = 1}

        \py{7 \% 2 = 1} $\rightarrow$ ข้าม

        \py{1 \% 2 = 1} $\rightarrow$ ข้าม

        \py{8 \% 2 = 0} $\rightarrow$ \py{c = 2}

        \py{3 \% 2 = 1} $\rightarrow$ ข้าม

        \highlight{แสดงผล \py{2}} (มีเลขคู่ 2 ตัว)
      \end{solbox}
  \end{columns}
\end{frame}

% ===========================================================================
%  SUMMARY
% ===========================================================================
\section{สรุป}

\begin{frame}{สรุปเนื้อหา}
  \begin{explanation}[สิ่งที่ได้เรียนรู้]
    \begin{columns}[T]
      \column{0.5\textwidth}
        \begin{itemize}
          \item \textbf{ชนิดข้อมูล}: int, float, bool, str
          \item \textbf{ตัวแปร}: \py{x = 5}, กฎการตั้งชื่อ
          \item \textbf{ตัวดำเนินการ}: \py{+ - * / // \% **}
          \item \textbf{เปรียบเทียบ}: \py{== != < > <= >=}
          \item \textbf{ตรรกะ}: \py{and or not}
          \item \textbf{ลำดับความสำคัญ}: วงเล็บ $\to$ \py{**} $\to$ \py{*/\%} $\to$ \py{+-}
        \end{itemize}
      \column{0.5\textwidth}
        \begin{itemize}
          \item \textbf{print()}: แสดงผลลัพธ์
          \item \textbf{if/elif/else}: เงื่อนไข + การเยื้อง
          \item \textbf{while}: วนซํ้าตามเงื่อนไข
          \item \textbf{break/continue}: ควบคุมการวนซํ้า
          \item \textbf{list}: \py{arr[i]}, \py{len(arr)}
          \item \textbf{string}: \py{s[i]}, \py{len(s)}, \py{s1 + s2}
        \end{itemize}
    \end{columns}
  \end{explanation}
  \begin{center}
    \large เอกสารนี้เป็นส่วนหนึ่งของเนื้อหาเตรียมสอบ \textbf{สอวน. คอมพิวเตอร์}
  \end{center}
\end{frame}

\end{document}
